\subsection{\texorpdfstring{\(\mathbf{R}^n\) 中的可和族}{R 的 n 次幂中的可和族}}

由于 \(\mathbf{R}^n\) 的每一点都具有可数的基本邻域系，一个族 \((\pmb{x}_i)\) （加法群 \(\mathbf{R}^n\) 中点的）只有当由指标 \(i\) （使得 \(\mathbf{x}_i \neq 0\) ）组成的集合可数时才是可和的（第 III 章 § 5 第 2 小节命题 1 的推论）；因此，本质上， \(R^{n}\) 中可和族的研究可以归结为可和序列的研究。尽管如此，出于与第 IV 章 § 7 中就 R 中的可和族所给过的同样理由，我们在下文中对指标集的基数将不施加任何限制。

命题 1. 点的族 \((\pmb{x}_t)_{t\in \mathbf{I}}\) （其中的点为 \(\pmb{x}_t = (x_{t,k})_{1\leq k\leq n}\) ，属于 \(\mathbf{R}^n\) ）是可和的，当且仅当 \(n\) 个实数族 \((x_{t,k})_{t\in \mathbf{I}}\) 中的每一个都在 \(\mathbf{R}\) 中可和。

这由第 III 章 § 5 第 4 小节命题 4 得出。

命题 1 的条件可以变换如下：

定理 1. 点的族 \((\pmb{x}_t)_{t\in \mathbf{I}}\) （其点属于 \(\mathbf{R}^n\) ）是可和的，当且仅当族 \((||\pmb{x}_t||)\) （即诸点 \(\pmb{x}_t\) 的欧氏范数所成的族）在 \(\mathbf{R}\) 中可和。

这可以毫不困难地由命题 1 、实数族的可和性条件（第 IV 章 § 7 第 2 小节定理 3）以及不等式

\[
\sup _ {1 \leqslant k \leqslant n} | x _ {t, k} | \leqslant | | x _ {t} | | \leqslant \sum_ {i = 1} ^ {n} | x _ {t, i} |,
\]

和比较原理（第 IV 章 § 7 第 1 小节定理 2）得出。

也可以稍许不同地进行，即先建立下列命题：

命题 2. 若 \((\pmb{x}_i)_{i\in \mathbf{I}}\) 是 \(\mathbf{R}^n\) 中点的任意一个有限族，则

\[
\sum_ {i \in \mathbf {I}} | | \boldsymbol {x} _ {i} | | \leqslant 2 n. \sup _ {\mathbf {J} \subset \mathbf {I}} \left| \left| \sum_ {i \in \mathbf {J}} \boldsymbol {x} _ {i} \right| \right|.\tag{1}
\]

因为若 \(\pmb{x}_i = (x_{ij})_{1\leqslant j\leqslant n}\) ，我们有 \(||\pmb{x}_i|| \leqslant \sum_{j=1}^{n} |x_{ij}|\) ，因此

\[
\sum_ {i \in \mathbf {I}} | | \boldsymbol {x} _ {i} | | \leqslant \sum_ {j = 1} ^ {n} \left(\sum_ {i \in \mathbf {I}} | x _ {i j} |\right).
\]

现在 \(\sum_{i\in \mathbf{I}}|x_{ij}| = \sum_{i\in \mathbf{I}}x_{ij}^{+} + \sum_{i\in \mathbf{I}}x_{ij}^{-}\) ，并且由于对 \(J\) （ \(I\) 的任意子集）我们有

\[
- \sum_ {i \in \mathbf {I}} x _ {i j} ^ {-} \leqslant - \sum_ {i \in \mathbf {J}} x _ {i j} ^ {-} \leqslant \sum_ {i \in \mathbf {J}} x _ {i j} ^ {+} \leqslant \sum_ {i \in \mathbf {I}} x _ {i j} ^ {+},
\]

由此得出

\[
\sum_ {i \in \mathbf {I}} | x _ {i j} | \leqslant 2. \sup _ {\mathbf {J} \subset \mathbf {I}} \left| \sum_ {i \in \mathbf {J}} x _ {i j} \right|.
\]

但是 \(\left|\sum_{i\in J}x_{ij}\right| \leqslant \left\| \sum_{i\in J}\boldsymbol{x}_i\right\|\) ，因此得到不等式 (1) 。

现在，定理 1 等价于下列命题（因为 \(\mathbf{R}^n\) 是完备群）：族 \((\pmb{x}_t)\) 满足柯西准则（第 III 章 § 5 第 2 小节定理 1），当且仅当族 \((||\pmb{x}_t||)\) 也满足柯西准则。而三角不等式表明这一条件是充分的，不等式 (1) 表明它是必要的。

此外，我们有不等式

\[
\left| \left| \sum_ {i} x _ {i} \right| \right| \leqslant \sum_ {i} | | x _ {i} | |,\tag{2}
\]

它是由关于有限部分和的类似不等式通过取极限得到的。

推论. 点的族 \((\pmb{x}_t)\) （属于 \(\mathbf{R}^n\) ）是可和的，当且仅当该族的有限部分和所成的集在 \(\mathbf{R}^n\) 中有界。

由定理 1 与三角不等式可知这一条件是必要的；而由不等式 (1) 与定理 1 可知它是充分的。

命题 3. 设 \((\pmb{x}_{\lambda})_{\lambda \in \mathbb{L}}\) 是 \(\mathbf{R}^m\) 中点的一个可和族， \((y_{\mu})_{\mu \in \mathbb{M}}\) 是 \(\mathbf{R}^n\) 中点的一个可和族，并设 \(f\) 是 \(\mathbf{R}^m \times \mathbf{R}^n\) 到 \(\mathbf{R}^p\) 中的一个双线性映射。那么族 \((f(x_{\lambda}, y_{\mu}))_{(\lambda, \mu) \in \mathbb{L} \times \mathbb{M}}\) 是可和的，并且我们有

\[
\sum_ {(\lambda , \mu) \in \mathbf {L} \times \mathbf {M}} f \left(\boldsymbol {x} _ {\lambda}, \boldsymbol {y} _ {\mu}\right) = f \left(\sum_ {\lambda \in \mathbf {L}} \boldsymbol {x} _ {\lambda}, \sum_ {\mu \in \mathbf {M}} \boldsymbol {y} _ {\mu}\right).\tag{3}
\]

为了证明族 \((f(x_{\lambda}, y_{\mu}))\) 是可和的，根据命题 1 ，只需证明由点 \(f(x_{\lambda}, y_{\mu})\) 在 \(R^{n}\) 中的坐标组成的 p 个族中的每一个都是可和的：换言之，我们可以限于 f 是双线性形式的情形；但对这样的形式 f 我们有

\[
f (\boldsymbol {x}, \boldsymbol {y}) = \sum_ {i, j} a _ {i j} x _ {i} y _ {j},
\]

因此我们被归结到 \(f(x, y) = x_i y_j\) 的情形，而在这种情形下结果已经证明（第 IV 章 § 7 第 3 小节命题 1）。

把函数 \(f\) 特殊化，我们特别地得到下列推论：

推论 1. 若 \((a_{\lambda})_{\lambda \in L}\) 是实数的一个可和族，且 \((x_{\mu})_{\mu \in M}\) 是 \(\mathbf{R}^n\) 中点的一个可和族，则族 \((a_{\lambda}x_{\mu})_{(\lambda, \mu) \in L \times M}\) 是可和的，且我们有

\begin{enumerate}
\def\labelenumi{(\arabic{enumi})}
\setcounter{enumi}{3}
\tightlist
\item
\end{enumerate}

\[
\sum_ {(\lambda , \mu) \in \mathbf {L} \times \mathbf {M}} a _ {\lambda} x _ {\mu} = \left(\sum_ {\lambda \in \mathbf {L}} a _ {\lambda}\right) \left(\sum_ {\mu \in \mathbf {M}} x _ {\mu}\right).
\]

推论 2. 若 \((\pmb{x}_{\lambda})_{\lambda \in \mathbb{L}}\) 与 \((\pmb{y}_{\mu})_{\mu \in \mathbb{M}}\) 是 \(\mathbf{R}^n\) 中点的两个可和族，则族 \((\pmb{x}_{\lambda}|\pmb{y}_{\mu})\) （参见第 VI 章 § 2 第 2 小节）在 \(\mathbf{R}\) 中可和，且我们有

\[
\sum_ {(\lambda , \mu) \in \mathbf {L} \times \mathbf {M}} (x _ {\lambda} | y _ {\mu}) = \left(\sum_ {\lambda \in \mathbf {L}} x _ {\lambda} \mid \sum_ {\mu \in \mathbf {M}} y _ {\mu}\right).\tag{5}
\]

